Es vol tancar un prat rectangular al costat d'un riu, de manera que el costat del riu no
necessita tanca. Es disposa de $400$ metres de tanca per als altres tres costats. Anomenem $y$
la longitud de cadascun dels dos costats perpendiculars al riu.

\begin{apartats}

\apartat[1,25]{1}
Comprova que l'àrea del prat és $A(y)=400y-2y^2$ i digues entre quins valors pot variar $y$.

\begin{solucio}
Si els dos costats perpendiculars al riu fan $y$ cadascun, el costat paral·lel al riu fa
$400-2y$. Per tant,
\[
A(y)=y\,(400-2y)=400y-2y^2 .
\]
Perquè el prat existeixi cal $y>0$ i $400-2y>0$: $y\in(0,200)$.
\end{solucio}

\apartat[1,25]{0,75}
Troba les dimensions del prat de més àrea i calcula aquesta àrea.

\begin{solucio}
$A'(y)=400-4y=0\iff y=100$.\\
Aleshores el costat paral·lel al riu fa $400-200=200$ metres: el prat fa
$\boxed{100\times200}$ metres i l'àrea màxima és $\boxed{20\,000}$ m$^2$.
\end{solucio}

\begin{nomesllarg}
\apartat{0,75}
Justifica que el valor trobat dona un màxim i no un mínim.

\begin{solucio}
$A''(y)=-4<0$ per a tot $y$: la funció és còncava i, per tant, en $y=100$ hi ha un
\textbf{màxim}.\\
També es veu pel signe de $A'$: positiva a $(0,100)$ i negativa a $(100,200)$.
\end{solucio}
\end{nomesllarg}

\end{apartats}
